\section{Conclusions}\label{sec:conclusions}

% --------------------- Sub Section 1 ---------------------

\begin{frame}{\secname}
	\framesubtitle{Wrapping up}

	Our macros bring \alert{stable specialization} to Rust, enabling more efficient code and improving readability and maintainability.

	\bigskip

	We are also the only approach to support specialization based on \alert{trait bounds}, \alert{lifetime bounds} and \alert{negative bounds}.

	\bigskip

	As an extension, we also support \alert{overloading}, a feature not natively available in Rust. % as it's seen as an extension of specialization

	\bigskip

	All of this works without compiler modifications, preserving \alert{soundness} and ensuring easy adoption by the Rust community.

\end{frame}

% --------------------- Sub Section 2 ---------------------

\begin{frame}{\secname}
	\framesubtitle{Final notes}

	\Large Interesting results:
	\normalsize
	\begin{itemize}
		\item We've written an article (\alert{Meta-Monomorphizing Specializations}), submitted to \alert{ECOOP}.
	\end{itemize}

	\bigskip

	\Large Interesting twist:
	\normalsize
	\begin{itemize}
		\item The approach can also be applied to other \alert{languages} % with powerful metaprogramming capabilities, needed to implement the macros
		      \\or even directly into the \alert{compiler}. % since the compiler already has all the necessary information during the compilation process and can simply replicate the macro's behavior natively
	\end{itemize}

\end{frame}

% --------------------- Sub Section 3 ---------------------

\begin{frame}
	\begin{center}
		\huge Thanks for your attention!
	\end{center}
\end{frame}